% ============================================================
% Question Bank: ch08-05 Stem-and-Leaf Plots
% Topic Title: Stem-and-Leaf Plots
% CCSS: Supplementary
% Questions: 27 (15 MC + 6 SA + 6 Graphical)
% ============================================================

% --- Multiple Choice Questions (15) ---

\begin{practiceQuestion}{8-5-q01}{mc}
\begin{questionText}
In a stem-and-leaf plot, the number $47$ would be displayed with a stem of:
\end{questionText}
\choiceA{$7$}
\choiceB{$47$}
\choiceC{$4$}
\choiceD{$0$}
\correctAnswer{C}
\explanation{In a stem-and-leaf plot, the stem is the tens digit and the leaf is the ones digit. For $47$, the stem is $4$ and the leaf is $7$.}
\end{practiceQuestion}

\begin{practiceQuestion}{8-5-q02}{mc}
\begin{questionText}
A stem-and-leaf plot shows the stem $6$ with leaves $1, 3, 5, 8$. What are the data values?
\end{questionText}
\choiceA{$6, 1, 3, 5, 8$}
\choiceB{$61, 63, 65, 68$}
\choiceC{$16, 36, 56, 86$}
\choiceD{$6.1, 6.3, 6.5, 6.8$}
\correctAnswer{B}
\explanation{Each leaf is combined with the stem to form a data value. Stem $6$ with leaves $1, 3, 5, 8$ gives $61, 63, 65, 68$.}
\end{practiceQuestion}

\begin{practiceQuestion}{8-5-q03}{mc}
\begin{questionText}
How many data values are in a stem-and-leaf plot that has stems $2, 3, 4, 5$ with $3, 5, 4, 2$ leaves respectively?
\end{questionText}
\choiceA{$4$}
\choiceB{$10$}
\choiceC{$14$}
\choiceD{$20$}
\correctAnswer{C}
\explanation{Count all the leaves: $3 + 5 + 4 + 2 = 14$. Each leaf represents one data value.}
\end{practiceQuestion}

\begin{practiceQuestion}{8-5-q04}{mc}
\begin{questionText}
A stem-and-leaf plot of test scores shows: stem $7$ with leaves $0, 2, 5, 5, 8$ and stem $8$ with leaves $1, 3, 6, 9$. What is the median of this data?
\end{questionText}
\choiceA{$75$}
\choiceB{$78$}
\choiceC{$80$}
\choiceD{$81$}
\correctAnswer{B}
\explanation{The $9$ ordered values are $70, 72, 75, 75, 78, 81, 83, 86, 89$. The middle (5th) value is $78$.}
\end{practiceQuestion}

\begin{practiceQuestion}{8-5-q05}{mc}
\begin{questionText}
Which of the following is an advantage of a stem-and-leaf plot over a histogram?
\end{questionText}
\choiceA{It works better for very large data sets}
\choiceB{It shows the exact data values while also showing the distribution shape}
\choiceC{It cannot display outliers}
\choiceD{It does not require ordered data}
\correctAnswer{B}
\explanation{A stem-and-leaf plot preserves the original data values (unlike a histogram, which groups data into bins). It also shows the shape of the distribution.}
\end{practiceQuestion}

\begin{practiceQuestion}{8-5-q06}{mc}
\begin{questionText}
A back-to-back stem-and-leaf plot is used to:
\end{questionText}
\choiceA{Display data for one group only}
\choiceB{Compare two data sets using a shared set of stems}
\choiceC{Show the mean of two groups}
\choiceD{Create a circle graph}
\correctAnswer{B}
\explanation{A back-to-back stem-and-leaf plot places two data sets on opposite sides of the same stems, making it easy to compare two groups.}
\end{practiceQuestion}

\begin{practiceQuestion}{8-5-q07}{mc}
\begin{questionText}
The stem-and-leaf plot shows daily temperatures (°F):

Stem $5$: $2, 5, 8$\\
Stem $6$: $0, 3, 3, 7$\\
Stem $7$: $1, 4$

What is the range of the data?
\end{questionText}
\choiceA{$22$}
\choiceB{$19$}
\choiceC{$25$}
\choiceD{$20$}
\correctAnswer{A}
\explanation{Minimum $= 52$, maximum $= 74$. Range $= 74 - 52 = 22°$F.}
\end{practiceQuestion}

\begin{practiceQuestion}{8-5-q08}{mc}
\begin{questionText}
A stem-and-leaf plot has the following data: $32, 35, 38, 41, 43, 45, 47, 50, 52$. What is the mode?
\end{questionText}
\choiceA{$35$}
\choiceB{$43$}
\choiceC{There is no mode}
\choiceD{$45$}
\correctAnswer{C}
\explanation{The mode is the most frequently occurring value. Since every value appears exactly once, there is no mode.}
\end{practiceQuestion}

\begin{practiceQuestion}{8-5-q09}{mc}
\begin{questionText}
The ages of participants in a workshop are shown in a stem-and-leaf plot:

Stem $2$: $1, 3, 5, 7, 9$\\
Stem $3$: $0, 2, 4, 6, 8$\\
Stem $4$: $1, 5$

What is the median age?
\end{questionText}
\choiceA{$30$}
\choiceB{$31$}
\choiceC{$32$}
\choiceD{$34$}
\correctAnswer{B}
\explanation{There are $12$ values. The median is the average of the 6th and 7th values. Ordered: ..., $29, 30, 32$, .... The 6th value is $30$ and the 7th is $32$. Median $= \frac{30+32}{2} = 31$.}
\end{practiceQuestion}

\begin{practiceQuestion}{8-5-q10}{mc}
\begin{questionText}
In a stem-and-leaf plot, the leaf unit is defined as $0.1$. The stem is $3$ and the leaf is $7$. What data value does this represent?
\end{questionText}
\choiceA{$37$}
\choiceB{$3.7$}
\choiceC{$0.37$}
\choiceD{$370$}
\correctAnswer{B}
\explanation{When the leaf unit is $0.1$, combine stem and leaf as $3.7$. The leaf unit tells you the place value of each leaf digit.}
\end{practiceQuestion}

\begin{practiceQuestion}{8-5-q11}{mc}
\begin{questionText}
A student creates a stem-and-leaf plot and accidentally puts the leaves out of order. What rule did the student break?
\end{questionText}
\choiceA{Stems must be in descending order}
\choiceB{Leaves must be listed in increasing order from left to right}
\choiceC{There must be exactly $5$ leaves per stem}
\choiceD{Stems and leaves must be the same digit}
\correctAnswer{B}
\explanation{In a properly constructed stem-and-leaf plot, leaves for each stem must be arranged in increasing order from left to right to show the distribution clearly.}
\end{practiceQuestion}

\begin{practiceQuestion}{8-5-q12}{mc}
\begin{questionText}
A stem-and-leaf plot shows:

Stem $1$: $8, 9$\\
Stem $2$: $0, 2, 4, 5, 5, 7$\\
Stem $3$: $1, 3$

How many values are in the $20$s?
\end{questionText}
\choiceA{$2$}
\choiceB{$4$}
\choiceC{$6$}
\choiceD{$10$}
\correctAnswer{C}
\explanation{The stem $2$ has leaves $0, 2, 4, 5, 5, 7$, representing $20, 22, 24, 25, 25, 27$. There are $6$ values in the $20$s.}
\end{practiceQuestion}

\begin{practiceQuestion}{8-5-q13}{mc}
\begin{questionText}
Data: $14, 17, 22, 25, 25, 28, 31, 36$. If you create a stem-and-leaf plot, stem $2$ would have which leaves?
\end{questionText}
\choiceA{$2, 5, 8$}
\choiceB{$2, 5, 5, 8$}
\choiceC{$14, 17$}
\choiceD{$22, 25, 25, 28$}
\correctAnswer{B}
\explanation{Values with tens digit $2$ are $22, 25, 25, 28$. Their leaves (ones digits) are $2, 5, 5, 8$.}
\end{practiceQuestion}

\begin{practiceQuestion}{8-5-q14}{mc}
\begin{questionText}
A stem-and-leaf plot displays quiz scores. The lowest score is $58$ and the highest is $97$. Which stems would appear in the plot?
\end{questionText}
\choiceA{$5, 6, 7, 8, 9$}
\choiceB{$58, 97$}
\choiceC{$0, 1, 2, ..., 9$}
\choiceD{$5, 9$}
\correctAnswer{A}
\explanation{The stems represent the tens digits. Since scores range from $58$ to $97$, the tens digits run from $5$ through $9$, so stems $5, 6, 7, 8, 9$ are needed.}
\end{practiceQuestion}

\begin{practiceQuestion}{8-5-q15}{mc}
\begin{questionText}
A stem-and-leaf plot has stems $4$ through $8$. No leaves appear next to stem $6$. What does this mean?
\end{questionText}
\choiceA{The stem was accidentally deleted}
\choiceB{No data values fall in the $60$--$69$ range}
\choiceC{All values are in the $60$s}
\choiceD{The data set is incomplete}
\correctAnswer{B}
\explanation{An empty stem means there are no data values with that tens digit. A stem of $6$ with no leaves means no values are between $60$ and $69$.}
\end{practiceQuestion}

% --- Short Answer Questions (6) ---

\begin{practiceQuestion}{8-5-q16}{sa}
\begin{questionText}
Data: $42, 45, 47, 53, 55, 58, 61, 64, 66, 69$. What is the median of this data set?
\end{questionText}
\correctAnswer{$56.5$}
\explanation{There are $10$ values. The median is the average of the 5th and 6th values: $\frac{55+58}{2} = \frac{113}{2} = 56.5$.}
\end{practiceQuestion}

\begin{practiceQuestion}{8-5-q17}{sa}
\begin{questionText}
A stem-and-leaf plot shows: stem $3$ with leaves $1, 4, 4, 7, 9$ and stem $4$ with leaves $0, 2, 6$. What is the mean of the data?
\end{questionText}
\correctAnswer{$37.875$}
\explanation{Values: $31, 34, 34, 37, 39, 40, 42, 46$. Sum $= 303$. Mean $= \frac{303}{8} = 37.875$.}
\end{practiceQuestion}

\begin{practiceQuestion}{8-5-q18}{sa}
\begin{questionText}
Create a stem-and-leaf plot (list the stems and their leaves) for this data: $25, 31, 28, 34, 22, 30, 27, 35, 29$.
\end{questionText}
\correctAnswer{Stem $2$: $2, 5, 7, 8, 9$; Stem $3$: $0, 1, 4, 5$}
\explanation{Sort the data: $22, 25, 27, 28, 29, 31, 30, 34, 35$. Stems are the tens digits. Stem $2$: leaves $2, 5, 7, 8, 9$. Stem $3$: leaves $0, 1, 4, 5$ (listed in order).}
\end{practiceQuestion}

\begin{practiceQuestion}{8-5-q19}{sa}
\begin{questionText}
A stem-and-leaf plot contains $15$ data values. The data has a range of $38$ and a minimum value of $24$. What is the maximum value?
\end{questionText}
\correctAnswer{$62$}
\explanation{Range $=$ max $-$ min, so max $= 24 + 38 = 62$.}
\end{practiceQuestion}

\begin{practiceQuestion}{8-5-q20}{sa}
\begin{questionText}
From the stem-and-leaf plot below, find the mode.

Stem $5$: $0, 3, 3, 7$\\
Stem $6$: $1, 3, 3, 3, 8$\\
Stem $7$: $2, 5$
\end{questionText}
\correctAnswer{$63$}
\explanation{The mode is the value that appears most often. The value $63$ appears $3$ times (three leaves of $3$ on stem $6$), which is more than any other value.}
\end{practiceQuestion}

\begin{practiceQuestion}{8-5-q21}{sa}
\begin{questionText}
A stem-and-leaf plot has stems $1$ through $5$. Stem $1$ has $2$ leaves, stem $2$ has $4$ leaves, stem $3$ has $6$ leaves, stem $4$ has $3$ leaves, and stem $5$ has $1$ leaf. How many data values are in the set, and which stem has the most values?
\end{questionText}
\correctAnswer{$16$ values; stem $3$ has the most.}
\explanation{Total: $2+4+6+3+1 = 16$. Stem $3$ has $6$ leaves, which is the most of any stem.}
\end{practiceQuestion}

% --- Graphical Questions (3 GMC + 3 GSA) ---

\begin{practiceQuestion}{8-5-q22}{gmc}
\begin{questionText}
Use the stem-and-leaf plot below to answer the question.

\begin{center}
\begin{tabular}{r|l}
\textbf{Stem} & \textbf{Leaf} \\
\hline
$3$ & $2 \quad 5 \quad 8$ \\
$4$ & $0 \quad 3 \quad 3 \quad 6 \quad 9$ \\
$5$ & $1 \quad 4 \quad 7$ \\
$6$ & $2 \quad 5$ \\
\end{tabular}
\end{center}

Key: $3 \mid 2 = 32$

What is the median of the data?
\end{questionText}
\choiceA{$43$}
\choiceB{$46$}
\choiceC{$44.5$}
\choiceD{$49$}
\correctAnswer{B}
\explanation{There are $13$ values in order: $32, 35, 38, 40, 43, 43, 46, 49, 51, 54, 57, 62, 65$. With $13$ values, the median is the 7th value, which is $46$.}
\end{practiceQuestion}

\begin{practiceQuestion}{8-5-q23}{gmc}
\begin{questionText}
The back-to-back stem-and-leaf plot below compares the ages of members of two book clubs.

\begin{center}
\begin{tabular}{r c l}
\textbf{Club A} & \textbf{Stem} & \textbf{Club B} \\
\hline
$9 \quad 7 \quad 5$ & $2$ & $1 \quad 3 \quad 6$ \\
$8 \quad 4 \quad 2 \quad 0$ & $3$ & $0 \quad 5 \quad 7$ \\
$5 \quad 1$ & $4$ & $2 \quad 4 \quad 8 \quad 9$ \\
$3$ & $5$ & $1 \quad 6$ \\
\end{tabular}
\end{center}

Key: $5 \mid 2 \mid 1 = $ Club A is $25$, Club B is $21$

Which club has more members who are $40$ or older?
\end{questionText}
\choiceA{Club A, with $3$ members}
\choiceB{Club B, with $6$ members}
\choiceC{They have the same number}
\choiceD{Club A, with $6$ members}
\correctAnswer{B}
\explanation{Club A ages $\geq 40$: stem $4$ ($41, 45$) and stem $5$ ($53$) $= 3$ members. Club B ages $\geq 40$: stem $4$ ($42, 44, 48, 49$) and stem $5$ ($51, 56$) $= 6$ members. Club B has more members aged $40$ or older.}
\end{practiceQuestion}

\begin{practiceQuestion}{8-5-q24}{gmc}
\begin{questionText}
Use the stem-and-leaf plot below showing the number of minutes students spent reading.

\begin{center}
\begin{tabular}{r|l}
\textbf{Stem} & \textbf{Leaf} \\
\hline
$1$ & $0 \quad 5 \quad 5$ \\
$2$ & $0 \quad 0 \quad 5 \quad 5 \quad 5$ \\
$3$ & $0 \quad 0 \quad 5$ \\
$4$ & $0$ \\
\end{tabular}
\end{center}

Key: $1 \mid 5 = 15$ minutes

What is the mode of the data?
\end{questionText}
\choiceA{$15$}
\choiceB{$20$}
\choiceC{$25$}
\choiceD{$30$}
\correctAnswer{C}
\explanation{Count each value: $10(1)$, $15(2)$, $20(2)$, $25(3)$, $30(2)$, $35(1)$, $40(1)$. The value $25$ appears $3$ times — more than any other — making it the mode.}
\end{practiceQuestion}

\begin{practiceQuestion}{8-5-q25}{gsa}
\begin{questionText}
Use the stem-and-leaf plot below showing heights (in inches) of plants in a garden.

\begin{center}
\begin{tabular}{r|l}
\textbf{Stem} & \textbf{Leaf} \\
\hline
$1$ & $2 \quad 4 \quad 6 \quad 8$ \\
$2$ & $0 \quad 3 \quad 5 \quad 7 \quad 9$ \\
$3$ & $1 \quad 4 \quad 6$ \\
\end{tabular}
\end{center}

Key: $1 \mid 2 = 12$ inches

Find the range and the median of the plant heights.
\end{questionText}
\correctAnswer{Range $= 24$ inches; median $= 24$ inches}
\explanation{Values in order: $12, 14, 16, 18, 20, 23, 25, 27, 29, 31, 34, 36$. Min $= 12$, max $= 36$; range $= 36 - 12 = 24$. There are $12$ values; median $= \frac{23 + 25}{2} = 24$ inches.}
\end{practiceQuestion}

\begin{practiceQuestion}{8-5-q26}{gsa}
\begin{questionText}
Create a stem-and-leaf plot for the following bowling scores: $112, 125, 118, 130, 127, 134, 121, 119, 132, 128$.

Show the stems and leaves in order.
\end{questionText}
\correctAnswer{Stem $11$: $2, 8, 9$; Stem $12$: $1, 5, 7, 8$; Stem $13$: $0, 2, 4$}
\explanation{Sort: $112, 118, 119, 121, 125, 127, 128, 130, 132, 134$. With leaf unit $= 1$ and stems as the first two digits: stem $11$: leaves $2, 8, 9$; stem $12$: leaves $1, 5, 7, 8$; stem $13$: leaves $0, 2, 4$.}
\end{practiceQuestion}

\begin{practiceQuestion}{8-5-q27}{gsa}
\begin{questionText}
The back-to-back stem-and-leaf plot compares scores on a science quiz for two classes.

\begin{center}
\begin{tabular}{r c l}
\textbf{Class X} & \textbf{Stem} & \textbf{Class Y} \\
\hline
$8 \quad 5 \quad 2$ & $6$ & $0 \quad 4 \quad 7$ \\
$9 \quad 6 \quad 3 \quad 0$ & $7$ & $2 \quad 5 \quad 8$ \\
$7 \quad 4 \quad 1$ & $8$ & $0 \quad 3 \quad 6 \quad 9$ \\
$5 \quad 2$ & $9$ & $1 \quad 4$ \\
\end{tabular}
\end{center}

Key: $2 \mid 6 \mid 0$ means Class X is $62$ and Class Y is $60$

Find the median score for each class.
\end{questionText}
\correctAnswer{Class X median $= 77.5$; Class Y median $= 79$}
\explanation{Class X sorted: $62, 65, 68, 70, 73, 76, 79, 81, 84, 87, 92, 95$ — $12$ values. Median $= \frac{76+79}{2} = 77.5$. Class Y sorted: $60, 64, 67, 72, 75, 78, 80, 83, 86, 89, 91, 94$ — $12$ values. Median $= \frac{78+80}{2} = 79$.}
\end{practiceQuestion}
